\documentclass[11pt,a4paper]{article}
\usepackage[utf8]{inputenc}
\usepackage[T1]{fontenc}
\usepackage{geometry}
\geometry{margin=1in}
\usepackage{hyperref}
\usepackage{enumitem}
\usepackage{booktabs}
\usepackage{longtable}
\usepackage{xcolor}
\usepackage{titlesec}
\usepackage{amsmath,amssymb}

\titleformat{\section}{\Large\bfseries\color{blue!70!black}}{}{0em}{}
\titleformat{\subsection}{\large\bfseries\color{blue!50!black}}{}{0em}{}
\titleformat{\subsubsection}{\normalsize\bfseries\color{blue!30!black}}{}{0em}{}

\title{\textbf{Replication Plan}\\[0.5em]
\large Quantitative Relationship Between Polarization Differences and the Zone-Averaged Shift Photocurrent\\[0.3em]
\normalsize OSTI ID: 1523841}
\author{Auto-generated Replication Blueprint}
\date{\today}

\begin{document}
\maketitle
\tableofcontents
\newpage

%---------------------------------------------------------------------
\section{Paper Overview}
%---------------------------------------------------------------------
This paper establishes a quantitative connection between the bulk photovoltaic (shift current) response and the difference in electronic polarization between valence and conduction bands. Using tight-binding models on 1D and 2D lattices (Rice--Mele model, Haldane model variants), the authors derive a zone-averaged shift-current expression relating the shift vector integral to polarization differences. The work is entirely analytic/numerical, with all model parameters specified.

%---------------------------------------------------------------------
\section{Detailed Workflow}
%---------------------------------------------------------------------

\subsection{Phase 1: Tight-Binding Model Construction}
\begin{enumerate}[label=\textbf{1.\arabic*}]
  \item \textbf{1D Rice--Mele model.}
    \begin{itemize}
      \item Two-site unit cell with alternating hopping $t_1, t_2$ and on-site energies $\pm \Delta$.
      \item Construct the $2 \times 2$ Bloch Hamiltonian $H(k)$.
    \end{itemize}
  \item \textbf{2D models} (Haldane-type, if applicable).
    \begin{itemize}
      \item Construct the Bloch Hamiltonian on the honeycomb lattice with next-nearest-neighbor hopping and complex phases.
    \end{itemize}
\end{enumerate}

\subsection{Phase 2: Band Structure and Berry Phase Calculations}
\begin{enumerate}[label=\textbf{2.\arabic*}]
  \item \textbf{Diagonalize $H(k)$} across the Brillouin zone on a dense $k$-mesh.
  \item \textbf{Compute Berry connection} $\mathcal{A}_n(k) = i \langle u_n(k) | \nabla_k | u_n(k) \rangle$ for each band $n$.
  \item \textbf{Compute Berry phase / polarization} $P_n = \frac{1}{2\pi} \oint \mathcal{A}_n(k) \, dk$ via the discretized Wilson loop.
  \item \textbf{Compute polarization difference} $\Delta P = P_{\text{cond}} - P_{\text{val}}$.
\end{enumerate}

\subsection{Phase 3: Shift Current Calculation}
\begin{enumerate}[label=\textbf{3.\arabic*}]
  \item \textbf{Compute the shift vector} $R_{cv}(k)$ for each $k$-point:
    \begin{equation*}
      R_{cv}(k) = \frac{\partial \phi_{cv}(k)}{\partial k} - [\mathcal{A}_c(k) - \mathcal{A}_v(k)]
    \end{equation*}
    where $\phi_{cv}$ is the phase of the interband matrix element.
  \item \textbf{Compute the shift current conductivity} $\sigma^{\text{shift}}(\omega)$ by integrating $R_{cv}(k)$ weighted by the joint density of states over the BZ.
  \item \textbf{Compute the zone-averaged shift current} and verify the relation to $\Delta P$.
\end{enumerate}

\subsection{Phase 4: Parameter Sweeps and Figure Reproduction}
\begin{enumerate}[label=\textbf{4.\arabic*}]
  \item \textbf{Sweep model parameters} ($t_1/t_2$, $\Delta$, etc.) as shown in the paper.
  \item \textbf{Plot shift current vs.\ polarization difference.}
  \item \textbf{Reproduce all figures and tables} from the paper.
  \item \textbf{Verify analytical limiting cases.}
\end{enumerate}

%---------------------------------------------------------------------
\section{Datasets}
%---------------------------------------------------------------------
No external datasets; all data is generated from model Hamiltonians.

\begin{longtable}{p{4.5cm} p{3.5cm} p{6cm}}
\toprule
\textbf{Dataset / Source} & \textbf{Type} & \textbf{Location / Notes} \\
\midrule
\endfirsthead
\toprule
\textbf{Dataset / Source} & \textbf{Type} & \textbf{Location / Notes} \\
\midrule
\endhead
Model parameters & Synthetic & Specified in paper text and figure captions \\
\midrule
Band structure data & Computed & Generated from Hamiltonian diagonalization \\
\midrule
Paper figures (reference) & Reference & For quantitative comparison of reproduced results \\
\bottomrule
\end{longtable}

%---------------------------------------------------------------------
\section{Models, Tools, and Software}
%---------------------------------------------------------------------

\subsection{Programming Languages}
\begin{itemize}
  \item \textbf{Python 3.8+} or \textbf{Julia 1.6+} --- for numerical implementation.
\end{itemize}

\subsection{Libraries and Frameworks}
\begin{longtable}{p{4.5cm} p{2.5cm} p{6.5cm}}
\toprule
\textbf{Tool / Library} & \textbf{Version} & \textbf{Purpose / URL} \\
\midrule
\endfirsthead
\toprule
\textbf{Tool / Library} & \textbf{Version} & \textbf{Purpose / URL} \\
\midrule
\endhead
NumPy & $\geq 1.21$ & Matrix diagonalization, array operations \\
\midrule
SciPy & $\geq 1.7$ & Numerical integration, interpolation \\
\midrule
Matplotlib & $\geq 3.4$ & Plotting band structures, shift current spectra \\
\midrule
PythTB & latest & Tight-binding model construction, Berry phase \newline \url{https://www.physics.rutgers.edu/pythtb/} \\
\midrule
Z2Pack & latest & Wilson loop and topological invariant calculations \newline \url{https://z2pack.greschd.ch/} \\
\midrule
WannierTools & latest & Optional: Wannier-based Berry phase calculations \newline \url{https://www.wanniertools.org/} \\
\bottomrule
\end{longtable}

\subsection{Hardware Requirements}
\begin{itemize}
  \item Any modern laptop; 1D models are computationally trivial; 2D models require a dense $k$-mesh but remain fast.
\end{itemize}

\end{document}
